\DocumentMetadata{
  pdfversion=2.0,
  pdfstandard=ua-2,
  lang=en-US,
  tagging=on,
  tagging-setup={math/setup=mathml-SE,tabsorder=structure}
}
\documentclass[11pt,letterpaper]{article}
\usepackage[margin=0.68in,includefoot]{geometry}
\usepackage{newcomputermodern}
\usepackage{amsmath,amscd,mathtools}
\usepackage{tikz,adjustbox}
\providecommand{\square}{\mdlgwhtsquare}
\usepackage[hidelinks]{hyperref}
\hypersetup{
  pdftitle={Algebra First-Year Exam, May 2019, Part 1},
  pdfauthor={University of Florida Department of Mathematics},
  pdfsubject={Graduate mathematics examination},
  pdfdisplaydoctitle=true,
  bookmarksopen=true
}
\setcounter{secnumdepth}{0}
\setlength{\parindent}{0pt}
\setlength{\parskip}{0.28em}
\setlength{\emergencystretch}{3em}
\setlength{\leftmargini}{2.15em}
\setlength{\leftmarginii}{2.65em}
\setlength{\leftmarginiii}{3em}
\renewcommand{\labelenumi}{\textbf{\theenumi.}}
\pagestyle{empty}
\raggedbottom
\allowdisplaybreaks
\newenvironment{examproblems}[1]{%
  \begin{enumerate}%
  \setcounter{enumi}{#1}%
  \setlength{\topsep}{0.35em}%
  \setlength{\partopsep}{0pt}%
  \setlength{\itemsep}{0.48em}%
  \setlength{\parsep}{0.18em}%
}{\end{enumerate}}
\newenvironment{examsubparts}{%
  \begin{enumerate}%
  \setlength{\topsep}{0.18em}%
  \setlength{\partopsep}{0pt}%
  \setlength{\itemsep}{0.16em}%
  \setlength{\parsep}{0.1em}%
}{\end{enumerate}}
\newenvironment{examinstructions}{%
  \par\begingroup\itshape
}{%
  \par\endgroup\vspace{0.18em}
}
\makeatletter
\renewcommand\section{\@startsection{section}{1}{0pt}%
  {-1.25ex plus -.25ex minus -.1ex}%
  {0.55ex plus .1ex}%
  {\normalfont\large\bfseries}}
\renewcommand{\@maketitle}{%
  \newpage\null\vskip -1em%
  {\footnotesize\bfseries
    \href{https://www.ufl.edu/}{University of Florida}, \href{https://math.ufl.edu/}{Department of Mathematics}\par}%
  \vskip -0.5em%
  \tagstructbegin{tag=Title}%
    {\LARGE\bfseries\href{https://gma.math.ufl.edu/exams/algebra-first-year/}{Algebra First-Year Exam}, May 2019, Part 1\par}%
  \tagstructend%
  \vskip 0.25em%
  \tagmcbegin{artifact}%
    \hrule height 1pt%
  \tagmcend%
  \vskip 1em%
}
\makeatother
\title{Algebra First-Year Exam, May 2019, Part 1}
\author{}
\date{}
\begin{document}
\maketitle
\thispagestyle{empty}
\begin{examinstructions}
Answer four problems. (If you turn in more, the first four will be graded.) Within reason, you may quote theorems as long as you state them clearly.
\end{examinstructions}
\begin{examproblems}{0}
\item
State and prove Sylow’s First Theorem.
\item
Prove that there is a non-abelian group of order \(2019\). (Hint: \(2019 = 3 \ast 673\) and \(673\) is a prime.)
\item
\(S_4\) is the symmetric group on four letters.
\begin{examsubparts}
\item[\textnormal{(a)}]
Prove that every simple subgroup of \(S_4\) is abelian.
\item[\textnormal{(b)}]
Use the result from (a) to prove that if~\(G\) is a non-abelian finite simple group, then every proper subgroup of~\(G\) has index at least~\(5\).
\end{examsubparts}
\item
Let~\(G\) be a finite group, and assume that~\(G\) acts on the set~\(\Omega\). For each \(\omega \in \Omega\), we denote by \(G_\omega\) the set of elements of~\(G\) that fix~\(\omega\).
\begin{examsubparts}
\item[\textnormal{(a)}]
Define what it means to say that the action is faithful.
\item[\textnormal{(b)}]
Define what it means to say that the action is transitive.
\item[\textnormal{(c)}]
Suppose~\(G\) acts on~\(\Omega\) faithfully and transitively. Suppose \(\nu \in \Omega\) and \(G_\nu\) is a cyclic group of order \(p^n\), where~\(p\) is a prime. Prove that there exists some \(g \in G\) such that \(G_\nu \cap G_{g\nu} = \{1\}\).
\end{examsubparts}
\item
Let~\(G\) be a group of order \(76 = 4 \cdot 19\).
\begin{examsubparts}
\item[\textnormal{(a)}]
Prove that~\(G\) contains a normal Sylow \(19\)-subgroup.
\item[\textnormal{(b)}]
Prove that the center of~\(G\) contains an element of order~\(2\).
\end{examsubparts}
\item
Let~\(G\) be a group, and let~\(M\) and~\(N\) be normal subgroups of~\(G\). Suppose that \(G/N\) and~\(N\) are simple groups. Suppose that \(NM \ne G\) and that \(N \ne M\). Prove that \(M = \{1\}\).
\end{examproblems}
\end{document}
